\documentclass{article}
\usepackage{decomp}
\begin{document}

\sourcefile{source/ta2023.pdf}
\decompnote{Publisher version of record (ACS Langmuir typeset PDF, 8 pages, open access under CC-BY 4.0). Left out: the running header ``Langmuir / pubs.acs.org/Langmuir / Article'' and the footer (page number, DOI link, ``Langmuir 2023, 39, 14396−14403'') on every page, the ``Open Access'' badge, the ACS ``Metrics \& More / Article Recommendations / Supporting Information'' link bar and the ``Read Online'' button on p.\,1, the vertical download stamp in the left margin of p.\,1 (``Downloaded via 188.216.158.91 on March 1, 2024 at 18:56:39 (UTC). See https://pubs.acs.org/sharingguidelines for options on how to legitimately share published articles.''), the ACS Publications logo and journal-cover thumbnail on p.\,1, and the ``Recommended by ACS'' advertisement box on p.\,8. The journal block (DOI line) is given once below.}

\pdfpage{1}
\noindent This article is licensed under CC-BY 4.0\par
\noindent pubs.acs.org/Langmuir\quad Article\par
\noindent https://doi.org/10.1021/acs.langmuir.3c01800\quad \textit{Langmuir} 2023, 39, 14396−14403\par

\section*{Atomistic Wear Mechanisms in Diamond: Effects of Surface Orientation, Stress, and Interaction with Adsorbed Molecules}

\noindent Huong T. T. Ta, Nam V. Tran,* and M. C. Righi*

\noindent\textbf{Cite This:} \textit{Langmuir} 2023, 39, 14396−14403

\begin{quote}
\textbf{ABSTRACT:} Despite its unrivaled hardness, diamond can be severely worn during the interaction with others, even softer materials. In this work, we calculate from first-principles the energy and forces necessary to induce the atomistic wear of diamond and compare them for different surface orientations and passivation by oxygen, hydrogen, and water fragments. The primary mechanism of wear is identified as the detachment of the carbon chains. This is particularly true for oxidized diamond and diamonds interacting with silica. A very interesting result concerns the role of stress, which reveals that compressive stresses can highly favor wear, making it even energetically favorable.
\end{quote}

%TOC
\begin{center}
\includegraphics[width=0.6\linewidth,keepaspectratio]{figures/toc.png}
\end{center}
\decompnote{The abstract graphic (table-of-contents graphic) is printed to the right of the abstract on p.\,1, without a caption or number; it is called ``toc'' in the file names.}
\alttext{Abstract graphic, a line plot with two y axes on a background shaded from green (top) through pale yellow to salmon/red (bottom). x axis ``Displacement (Å)'', linear, ticks 0 to 2.5 in steps of 0.5. Left y axis (blue) ``Energy (eV)'', −1 to 10; right y axis (red) ``Forcce (nN)'' (so spelled), 0 to 12 in steps of 2 with minor ticks below 0. Blue circles joined by a smooth blue line (energy): about 0 at 0 Å, 0.33 at 0.25, 1.47 at 0.5, 1.98 at 0.75, 2.6 at 1.0, a plateau of about 2.45 from 1.25 to 1.75, 2.7 at 2.0, a maximum of about 3.3 at 2.25 and 3.06 at 2.5. Red squares joined by a dashed red line (force): about 0 at 0, 4.2 at 0.25, 7.65 at 0.5, 2.8 at 0.75, 5.35 at 1.0, −0.27 at 1.25, −0.06 at 1.5, 0.6 at 1.75, 2.6 at 2.0, 4.7 at 2.25 and 0.15 at 2.5 nN; the dashed curve dips below zero between 1.25 and 1.5. Five round, white-vignetted structure insets (ball-and-stick side views of the C(110) slab, C black and grey, O red) sit above the markers at about 0.25, 0.5, 0.75, 1.0 and 2.25 Å, showing the O atom pulled up and a chain of C atoms growing beneath it (a vertical chain of several C atoms in the last one). No caption.}
\figuredata{(single panel)}{figures/toc.csv}{figures/toc.py}{tier 3, digitized from the embedded 300 ppi raster image (no vector data in the PDF); marker centres located by colour masks and axis fits on the major ticks; covered markers read from their visible edges or from the curve on both sides (flagged in the CSV); see figures/toc\_extract.py}

\noindent\textbf{Received:} June 29, 2023\\
\textbf{Revised:} September 10, 2023\\
\textbf{Published:} September 27, 2023

\noindent © 2023 The Authors. Published by American Chemical Society
\decompnote{The dates (lower right) and the copyright line (bottom centre) are printed at the foot of p.\,1; they are placed here after the abstract.}

\section*{1. INTRODUCTION}

The excellent wear resistance and ultralow friction make diamond an ideal choice for various tribological applications such as cutting tools, bearings, micro-electromechanical systems, and coatings.\textsuperscript{1−3} However, diamond is not completely inert to wear, especially under extreme conditions which experience high loads and a severe environment. The impact of the environment and/or the extreme working conditions result in material loss, high friction, and surface degradation and seriously undermine the functionality of the diamond-based devices.\textsuperscript{4−7} Understanding the wear mechanism on diamond surfaces is of utmost significance, not only for enhancing its applications but also for advancing material synthesis.

Several studies have described the chemical modifications of diamond surfaces under operating conditions, such as humid air, where atmospheric gases (H\textsubscript{2}, O\textsubscript{2}, and H\textsubscript{2}O) are present. The presence of H\textsubscript{2} or H\textsubscript{2}O can create surface hydrogenation or hydroxylation, which have been shown to be beneficial for reducing friction and wear at the contacting interfaces.\textsuperscript{8−11} The primary mechanism of reduced friction is the surface passivation when carbon atoms from dangling bonds at the surface are terminated by H/OH. Recently, it has been shown that O\textsubscript{2} dissociation is even kinetically and thermodynamically more favorable on diamond surfaces than H\textsubscript{2} and H\textsubscript{2}O,\textsuperscript{12} making surface oxidation easier than hydrogenation or hydroxylation. The dissociative adsorption O\textsubscript{2} not only alters diamond surface morphology but also modifies its electronic structure such as narrowing the band gap,\textsuperscript{13} changing the surface reactivity,\textsuperscript{14} or conductivity.\textsuperscript{15} In particular, O\textsubscript{2} dissociation induces the breaking of surface dimer bonds, leading to the de-reconstruction on the C(100) surface.\textsuperscript{16} It has been found that the presence of oxidizing agents can promote the desorption of CO and CO\textsubscript{2} molecules,\textsuperscript{17,18} which is the principle of oxidative etching and fabrication of diamond.\textsuperscript{13}

Great effort has been devoted to understanding the atomic mechanisms of chemical and mechanical polishing processes of diamond surfaces,\textsuperscript{18−21} and several mechanisms have been proposed. In particular, the sp\textsuperscript{3}-to-sp\textsuperscript{2} transformation has been proposed to explain the experimental observation of amorphous carbon and amorphous wear particles.\textsuperscript{22} By MD simulations, it was shown that C−C bond breaking on the C(110) surface starts at the weak C−C bonds which connect the C−C zigzag chain to the bulk diamond when sliding against silica. In this case, the C−C bond breaking can be initiated through strong Si−C and O−C bonds where Si and O play as mechanically supported pilot atoms.\textsuperscript{23,24} This event was not observed when sliding against silicon,\textsuperscript{24} highlighting the role of counter surfaces on the diamond polishing process. By a tight-binding quantum chemical molecular dynamics simulation, Kubo et al. proposed two mechanisms of wear, i.e., atom-by-atom and sheet-by-sheet removal of diamond polishing in the presence of OH radicals acting as an oxidizing agent.\textsuperscript{21} In H\textsubscript{2}O\textsubscript{2} solution, it has been reported that carbon atoms can react with decomposed −H, −OH, and −O, leading \pdfpage{2}to the formation of C−H, C−OH, and C−O. By MD simulations, the authors found that C atoms can be removed in the forms of CO, CO\textsubscript{2}, or carbon chains as a result of the combination of chemical and mechanical effects.\textsuperscript{18} Another aspect is the dependency of wear on surface orientations.\textsuperscript{20,22,25} While C(111) is reported as the hardest plane to be polished,\textsuperscript{22} C(001) and C(110) are more vulnerable and easier to be worn.\textsuperscript{25} This point has been proven and clearly shown on clean surfaces. Nevertheless, the effects of oxidation and the adsorption of atmospheric species on diamond polishing remain unclear. A detailed investigation of the effect of adsorbed O and other chemical groups on diamond polishing is of significance for its widespread applications in practical conditions.

In this study, we present a new computational procedure grounded in first-principles calculations to systematically characterize different atomistic wear mechanisms and quantify their associated energy costs and forces. By applying this methodology, we investigate the influence of surface orientation, namely, the C(001), C(110), and C(111) surfaces of diamond, as well as the adsorption of O, H, and OH species, interaction with silica, and surface stress (compression and expansion) on the wear mechanism. Our findings reveal that the adsorption of an oxygen (O) atom in the ketone configuration facilitates the detachment of carbon atoms from the diamond surface, thus initiating the wear process. Additionally, we uncover a wear mechanism on the diamond surface involving the formation of extended carbon chains. The C(110) diamond surface was identified as the surface that was most susceptible to wear in our study. Conversely, the adsorption of other species has a limited impact on the wear mechanism. Furthermore, we demonstrate that the initiation of carbon atom wear on the diamond surface necessitates the presence of a robust bond from the molecule or counter surface, such as a double bond or two single bonds. These insights shed light on the underlying mechanisms governing wear phenomena in the context of diamond surfaces and offer valuable implications for understanding and controlling wear processes at the atomic level.

\section*{2. SIMULATION METHODS}

In this research, density functional theory (DFT) simulations were performed using the Quantum ESPRESSO software.\textsuperscript{26} The exchange-correlation term was described using the generalized gradient approximation (GGA) as parametrized by Perdew−Burke−Ernzerhof.\textsuperscript{27} To account for the long-range van der Waals interactions, a semiempirical correction by Grimme (D2) was used.\textsuperscript{28,29} This combination of GGA-PBE and D2 has been found to provide an optimal balance between accuracy and computational cost, as indicated by previous studies.\textsuperscript{30,31} The convergence threshold was set at 10\textsuperscript{−4} Ry for the total energy and 10\textsuperscript{−3} Ry/bohr for the ionic forces. The self-consistent electronic (SCF) loop was set to converge at 10\textsuperscript{−6} Ry. Spin polarization was included in all calculations, as the presence of surface dangling bonds and dissociated molecules could result in magnetization in the system.

In this study, three distinct diamond surfaces were evaluated, namely, the C(110), the dimer-reconstructed C(001) and the Pandey-reconstructed C(111) surfaces. The simulations utilized a large orthorhombic supercell with a minimum lateral dimension of 8.74 Å in order to minimize interaction with periodic replicas. The C(110) and C(001) surfaces were modeled by using a 4 × 4 in-plane supercell, while the C(111) surface was simulated by using a 4 × 3 in-plane slab. The slab thickness used to model the C(110) corresponds to 7 atomic layers. Meanwhile, slabs of 8 and 10 atomic layers have been used for the C(111) and C(001) surfaces, respectively. The thickness values were selected based on previous studies,\textsuperscript{32−35} for obtaining accurate structural and energetic properties. A 20 Å vacuum region was included in the supercell to separate each slab from its periodic replica along the [001] direction. All simulations employed a plane wave cutoff of 30 Ry and a Monkhorst−Pack (MP) k-point mesh of 2 × 2 × 1. The cutoff energy, k-point mesh, and vacuum thickness were tested to ensure that the energy error was less than 3 meV/atom.

\section*{3. RESULTS AND DISCUSSION}

\noindent\textbf{3.1. Wear of O-Containing Diamond Surfaces.} We first investigate the effect of oxygen adsorption on the wear mechanism of diamond surfaces including C(110), C(001), and C(111)-Pandey compared with that of the clean ones. The initial structures for the calculations are presented in Figure 1.

%FIG1
\begin{center}
\includegraphics[width=0.75\linewidth,height=0.55\textheight,keepaspectratio]{figures/fig1.png}
\end{center}
\noindent\textbf{Figure 1.} Initial structures for clean diamond surfaces and the dissociative adsorption of oxygen on the three diamond surfaces. C (gray) and O (red). The bold color is to show the carbon atoms at the top layer. The green boxes show the constrained atoms.
\alttext{A 3 × 3 grid of ball-and-stick top views of diamond slabs, in light-blue frames. Columns are labelled at the top ``C(110)'', ``C(001)'', ``C(111)''; rows are labelled at the left ``Clean'', ``Ether'', ``Ketone''. Carbon atoms are grey, those of the top layer drawn darker (black); oxygen atoms red. Clean row: C(110) shows zigzag rows of dark top-layer atoms over a hexagonal grey network; C(001) shows a square grid with pairs of dark (dimer) atoms in columns at the left and right; C(111) shows a hexagonal network with a dark zigzag row. Ether row: one or two red O atoms bridge pairs of surface C atoms (drawn on the C−C bond). Ketone row: a red O atom sits on top of a single surface C atom; on C(110) a second O bridges a bond at the upper right, on C(001) a second O is at the left, on C(111) a second O bridges a bond at the right. In every panel a small green square box marks one atom (the constrained atom per the caption): the central top-layer C atom in the clean row and the O atom (the bridging O at the right in the ether row, the ketone O in the ketone row). No axes or data.}

The dissociative adsorptions of the oxygen molecule are considered with both ether (oxygen atoms form C−O−C bonds with the carbon atoms of diamond surfaces) and ketone (oxygen atoms form a double bond with a single carbon atom of diamond surfaces C═O) configurations (Figure 1). Previous works showed that ketone configurations are more stable for the C(110) and C(001) at low adsorption coverage.\textsuperscript{13} Meanwhile, the configuration is highly unfavorable for the C(111) reconstructed surface.\textsuperscript{12}

In practical working conditions, diamond surfaces are in contact with a counter body. The extreme condition results in the bond formation between the oxygen or unpassivated carbon and other atoms of the counter surface.\textsuperscript{9,36} The rubbing of the two surfaces leads to the stretching of the interfacial bonds and lifting of the oxygen atom or the underlying carbon off the surface. In an ideal scenario and to compare the wear mechanisms induced on the clean and different passivation of diamond surfaces, we mimic the effect of wear by pulling one atom (the adsorbed O atom or the C atom for the case of clean surface) along the vertical direction. In particular, for each step, \pdfpage{3}the oxygen atom is displaced by 0.25 Å, followed by an atomic relaxation with the position of the O atom fixed along the \textit{z} direction. The pulling energy is calculated by subtracting to the energy of the system where one atom is displaced out of its equilibrium position the energy at equilibrium. Meanwhile, the pulling force is the residual force along the vertical direction on the constrained atom, which is calculated by the Quantum Espresso program. In Figure 2a the energy increase due to the displacement of the pulled atom is reported as a function of the atom−surface distance (the corresponding forces are reported in Figure S1). The arrows indicate the first atom’s detachment from the surface. We show that very high pulling energies are required to detach the carbon atom from the clean surface. In particular, pulling energies of 2.95, 4.73, and 8.54 eV are required to remove the first carbon atom from the C(110), C(001), and C(111) clean surfaces, respectively. The high extreme pulling energy of C(111) is consistent with previous experimental and theoretical studies indicating the surface is very stable.\textsuperscript{12} On the other hand, the lower stability of C(110), which has higher surface energy than the other two surfaces, explains the reason for the lower energy cost for the wear processes. Another interesting observation common to all the considered surfaces is that the removal of carbon atoms is likely in the form a carbon chain, which is consistent with what found in diamond polishing.\textsuperscript{18} It is worth mentioning that in the counter surface environment, the carbon chain with carbon dangling bonds can react with other species, promoting a permanent loss of carbon atoms of the surface. Thus, the carbon chain established from our DFT calculations can be considered as the precursor for more complex tribochemical reactions with counter materials or environmental species.

%FIG2
\begin{center}
\includegraphics[width=0.95\linewidth,height=0.55\textheight,keepaspectratio]{figures/fig2.png}
\end{center}
\noindent\textbf{Figure 2.} Wear of clean (a), ether (b), and ketone (c) containing diamond surfaces.
\decompnote{Figure 2 is printed across the top of p.\,3, above the first text line of that page; it is placed after the paragraph that runs from p.\,2 onto p.\,3.}
\alttext{Three line plots side by side, each with x axis ``Displacement (Å)'' and y axis ``Energy (eV)'' (printed once, at the left) from −1 to 10 with major ticks every 1 eV; markers joined by smooth curves; legend (in a) green squares ``C(001)'', red triangles ``C(110)'', blue circles ``C(111)''. Each panel has colour-framed ball-and-stick structure insets (side views; frame colour = series), a blue numbered call-out ``1'' on a dotted vertical line at 0.25 Å and a call-out ``2'' on a dotted line at the right end, and short coloured arrows marking points discussed in the text.
(a) ``(a) Clean surface'', x 0 to 2.5 Å. C(001): 0, 0.41, 1.48, 3.06, a peak of 4.74 eV at 1.0 Å (green arrow), then 3.72, 3.49, 4.0, 5.0 and 6.45 eV; its last visible markers sit at about 1.8, 2.05 and 2.3 Å, not on the 0.25 Å grid, and the curve continues under the green-framed inset. C(110): a dip to about −0.14 eV at 0.25 Å, then 1.1, 1.4, 2.1, 2.93 eV at 1.25 Å (red arrow), 3.1, 3.48, 4.21 at 2.0 Å, 3.80 and 3.82 eV. C(111): rises steeply to 4.13 at 1.0, 5.58, 6.97, a local maximum of 8.52 eV at 1.75 Å (blue arrow), 7.89 at 2.0, 8.87 and 9.10 eV at 2.5 Å. Left insets (top to bottom, blue, green, red frames): the three clean slabs; right insets: blue frame (top) a slab with a chain of C atoms pulled upward, green frame a slab with a short C chain, red frame (lower right) a slab with a short C chain on a disordered surface.
(b) ``(b) Ether'', x 0 to 2 Å; the red and blue markers are pale (semi-transparent). All three curves rise from 0 to a plateau: C(001) to a maximum of 5.01 eV at 1.25 Å (green arrow) and about 4.6–4.9 eV after it; C(110) to about 4.3–4.6 eV (red arrow near 1.25 Å); C(111) to about 3.5 eV at 1.0 Å (blue arrow) and 3.7–4.0 eV after it. Left insets: ether-configured slabs with O (red) bridging two C atoms; right insets: the O atom detached above an intact surface.
(c) ``(c) Ketone'', x 0 to 2.5 Å. C(001): 0.3 at 0.25, 1.28 at 0.5, 2.76 eV at 0.75 Å (green arrow), then a flat plateau of about 2.47–2.54 eV from 1.0 to 2.5 Å. C(110): 1.48, 1.96 eV at 0.75 Å (red arrow), 2.58 at 1.0, about 2.4 from 1.25 to 1.75, then 2.70, 3.26 and 3.03 eV. C(111): slow rise (0.13, 0.48, 1.02, 1.89), a local maximum of 3.13 eV at 1.25 Å (blue arrow), 2.50 at 1.5, then 3.48, 4.84, a maximum of 6.17 eV at 2.25 Å and 4.18 eV at 2.5 Å. Left insets: ketone slabs (O atop single C atoms); right insets: blue frame a C chain with O on top standing up from the surface, green frame a detached C−O pair above the surface, red frame (lower right) a chain of C atoms with O at its end rising from the surface.}
\figuredata{a}{figures/fig2a.csv}{figures/fig2a.py}{tier 3, digitized from the embedded 300 ppi raster image (no vector data in the PDF); marker centres located by colour masks and axis fits on the major ticks; covered markers read from their visible edges or from the curve on both sides (flagged in the CSV); see figures/fig2a\_extract.py}
\figuredata{b}{figures/fig2b.csv}{figures/fig2b.py}{tier 3, digitized from the embedded 300 ppi raster image (no vector data in the PDF); marker centres located by colour masks and axis fits on the major ticks; covered markers read from their visible edges or from the curve on both sides (flagged in the CSV); see figures/fig2b\_extract.py}
\figuredata{c}{figures/fig2c.csv}{figures/fig2c.py}{tier 3, digitized from the embedded 300 ppi raster image (no vector data in the PDF); marker centres located by colour masks and axis fits on the major ticks; covered markers read from their visible edges or from the curve on both sides (flagged in the CSV); see figures/fig2c\_extract.py}

The adsorption of oxygen in the ether configuration is unlikely to cause wear of the diamond surfaces through the detachment of the oxygen atoms, which occurs without causing any deformation of the underlying carbon−carbon bonds (Figure 2b). The energy trend is similar for all three surfaces, in which the system energy increases as the oxygen atom is pulled up and becomes flat when the oxygen atom is detached from the surface. On the contrary, the pulling up of oxygen atoms in the carbonyl (ketone) group can lead to the detachment of carbon atoms from the surfaces. We found that the wear mechanism depends on the surface orientation. For the C(110) and C(111) surfaces that terminate with zigzag chains, the displacement of oxygen atoms can cause progressive detachment of the surface chains, which is consistent with experimental findings on wear regime of diamond.\textsuperscript{23,37} Both the force (Figure S1) and the energy (Figure 2c) are much higher for the C(111) surface, as this surface is highly stable.\textsuperscript{12} Meanwhile, the wear on the C(001) surface will likely form a CO molecule rather than a carbon chain (Figure 2c). Compared to the case of the clean surface, the presence of oxygen atom in the ketone configuration helps to reduce the energy costs for carbon detachment. In particular, the pulling energies required to detach the first carbon atom from the surface are reduced to 1.97, 2.51, and 3.13 eV for the C(110), C(001), and C(111) surfaces. Among the three surface directions, the C(110) is the most easily wearable because the associated energy cost and restoring forces (Figure S1) are lower than for the other surfaces. It is worth noting that when the diamond surface is in contact with a counter surface, the pulling process proceeds through an interfacial bond connecting the two surfaces. Thus, the O atom of the C═O bond is bonded to other atoms, and a double C═O bond cannot be established. However, the strength of the C═O bond can be resembled when the carbon atom is bonded to two single bonds, as in the bidentate configuration of the silicate. More details about this adsorbed configuration are discussed in Section 3.2.

As our calculation shows that C(110) is the easiest surface orientation to be polished, further calculations of wear mechanisms and effects of adsorbed species will focus on this surface. First, the energy changes and detailed atomistic mechanism of carbon detachment from the C(110) surface are depicted in Figure 3. When the oxygen atom is displaced from the surface (step 1 → 2), the two C−C bonds are stressed, \pdfpage{4}causing an increase in the energy and restoring force acting on the displaced atom. After a C−C bond is broken (indicated by the green arrow), the force decreases to a lower value. If the displacement of the O atom is continued, the restoring force on the O atom increases again (step 3 → 4) until the second C−C bond with the surface is broken. The dissociation of one of the two C−C bonds is a prerequisite for the establishment of a carbon chain. Primarily, it facilitates upward traction of the chain and detachment of the carbon atom from the substrate. Second, the cleavage of one C−C bond results in the transformation of the remaining C−C bond into a double bond. The presence of these double bonds is essential to enduring the mechanical forces exerted during the pulling process. As this mechanism is repeated, the carbon chain becomes longer and longer. It is also worth mentioning that the longer the carbon chain, the longer the displacement that it takes for the force to increase. The result is due to the elastic properties of the carbon chain that as a collection of springs added in series becomes less stiff as a new C−C bond is incorporated into the chain.

%FIG3
\begin{center}
\includegraphics[width=0.7\linewidth,height=0.55\textheight,keepaspectratio]{figures/fig3.png}
\end{center}
\noindent\textbf{Figure 3.} Energy (blue) and force (red) evolution under detachment of a carbon chain on the C(110) surface during the wear process.
\decompnote{Figure 3 is printed at the bottom of the right column of p.\,3, and Figures 4 and 5 across the top of p.\,4; all three are placed after the paragraph that runs from p.\,3 onto p.\,4, so the \textup{[p.\,4]} marker above comes before them.}
\alttext{Line plot with two y axes. x axis ``Displacement (Å)'', 0 to 2.5. Left y axis in blue ``Energy (eV)'', −1 to 10; right y axis in red ``Forcce (nN)'' (so spelled), 0 to 12 in steps of 2. Blue circles on a smooth solid blue curve (energy): about 0, 0.34, 1.47, 1.98, 2.61 (1.0 Å), 2.47, 2.44, 2.47, 2.70, 3.28 (2.25 Å, maximum) and 3.06 eV (2.5 Å). Red squares on a dashed red curve (force): about 0.0, 4.2, 7.65 (0.5 Å, maximum), 2.8, 5.35 (1.0 Å), −0.25, −0.05, 0.6, 2.6, 4.74 (2.25 Å) and 0.17 nN (2.5 Å); the dashed curve dips to about −0.9 nN near 1.35 Å. Black dotted vertical lines with blue numbered call-outs 1–6 mark 0.25, 0.5, 0.75, 1.0, 2.25 and 2.5 Å. Along the top, six light-blue-framed side-view insets numbered 1–6 show the C(110) slab (C black/grey, O red): in 1 and 2 a label ``C-C'' with a green arrow points at a C−C bond of the lifted ketone C atom; from 3 to 6 the O atom stands on top of a vertical chain of C atoms pulled out of the surface, the chain getting longer (about four C atoms, tilted, in inset 6).}
\figuredata{(single panel)}{figures/fig3.csv}{figures/fig3.py}{tier 3, digitized from the embedded 300 ppi raster image (no vector data in the PDF); marker centres located by colour masks and axis fits on the major ticks; covered markers read from their visible edges or from the curve on both sides (flagged in the CSV); see figures/fig3\_extract.py}

%FIG4
\begin{center}
\includegraphics[width=0.85\linewidth,height=0.55\textheight,keepaspectratio]{figures/fig4.png}
\end{center}
\noindent\textbf{Figure 4.} Top and lateral views of the O adatom arrangements considered for the C(110) surface (a). Energy change during the out-of-plane displacement of the O atom belonging to the ketone group (b).
\alttext{(a) Four ball-and-stick panels in light-blue frames, columns labelled ``O3'' and ``O4'': top views (upper row) and side views (lower row) of the C(110) slab with C black (top layer) and grey and O red. Both have the ketone O on a lifted C atom plus bridging O atoms; with O3 the extra O is farther from the ketone site, with O4 an O sits on the C atom next to it (side view: a ketone O above the centre and further O atoms beside it).
(b) Line plot, x ``Displacement (Å)'' 0 to 2.5, y ``Energy (eV)'' −1 to 7; legend O3 red squares, O4 blue circles (both semi-transparent). O3: 0, 0.34, 1.26, 1.49, 2.05, 2.78, a local maximum 3.16 eV at 1.5 Å, 2.43 at 1.75, 3.03, 4.11 and 5.31 eV at 2.5 Å. O4: 0, 0.38, 1.33, 2.53, 3.72 at 1.0 Å, then flat at about 3.76–3.82 eV to 2.5 Å. Call-out 1 on a dotted line at 0 Å, call-out 2 on a dotted line at 2.5 Å. Insets: left blue-framed and red-framed initial side views; top right blue frame: a C−O pair detached above the surface (O4 end state); lower right red frame: a vertical chain of C atoms with O on top rising from the surface (O3 end state).}
\figuredata{b}{figures/fig4b.csv}{figures/fig4b.py}{tier 3, digitized from the embedded 300 ppi raster image (no vector data in the PDF); marker centres located by colour masks and axis fits on the major ticks; covered markers read from their visible edges or from the curve on both sides (flagged in the CSV); see figures/fig4b\_extract.py}

%FIG5
\begin{center}
\includegraphics[width=0.85\linewidth,height=0.55\textheight,keepaspectratio]{figures/fig5.png}
\end{center}
\noindent\textbf{Figure 5.} Top and lateral views of the chemical groups adsorbed on the C(110) surface (a). Energy during the out-of-plane displacement of the H, OH, and O groups (b). C (gray), O (red), and H (white).
\alttext{(a) Eight ball-and-stick panels in light-blue frames, columns labelled ``OH'', ``OOH'', ``H\textsubscript{2}'', ``H\textsubscript{2}O'': top views (upper row) and side views (lower row) of C(110) with the adsorbed groups (C black/grey, O red, H white/pink). OH: an O atom on a surface C and an H atom nearby; OOH: an O−O−H group and an O atom; H\textsubscript{2}: two H atoms on surface C atoms; H\textsubscript{2}O: an OH group and an H atom.
(b) Line plot, x ``Displacement (Å)'' (the 0 Å tick label is hidden under an inset; visible labels 0.5 to 2.5), y axis −1 to 7 without an axis title; legend OH (black line, grey square), H2 blue circles, H2O red squares, OOH green triangles; the OH data are drawn as small grey diamonds. OH: 0, 0.3, 1.35, 2.66 at 0.75 Å, a flat stretch of about 2.6–2.7 eV to 1.75 Å, then 2.80, 3.34 and 4.19 eV at 2.5 Å (black curve, call-out 4 at 2.5 Å). H2: sigmoidal rise 0, 0.06, 0.53, 1.58, 2.70, 3.83, 4.62 (1.5 Å, call-out 2), 4.83, 4.95 eV at 2.0 Å. H2O: just below H2: 0.11, 0.54, 1.38, 2.53, 3.58, 4.23, 4.50, 4.63 eV at 2.0 Å. OOH: 0.30, 0.91, 1.96, 3.18, 4.11, 4.55 eV at 1.5 Å, where it ends. Dotted vertical lines at 0, 1.5, 2.0 and 2.5 Å with call-outs 1, 2, 3 and 4 (a second ``3'' at the upper right). Insets: left column blue-, red- and green-framed initial side views, a black-framed one at the bottom left (OH, an O on a C atom); top centre green frame: an O−H fragment detached above the surface; upper right blue frame: an H atom above the surface; red frame: an O−H fragment above the surface; black frame (right): a chain of C atoms with O on top standing up from the surface.}
\figuredata{b}{figures/fig5b.csv}{figures/fig5b.py}{tier 3, digitized from the embedded 300 ppi raster image (no vector data in the PDF); marker centres located by colour masks and axis fits on the major ticks; covered markers read from their visible edges or from the curve on both sides (flagged in the CSV); see figures/fig5b\_extract.py}

We examined the effects of the presence of other O atoms adsorbed in the surrounding area of the ketone group (Figures 4a and S2). The previous calculation shows that only ketone configuration could facilitate the formation of carbon chain. Consequently, we initiated our investigation by considering the C(110) structure, which incorporates the ketone configuration, as depicted in Figure 1. Subsequently, we introduce additional oxygen atoms one by one close to the ketone site, which is chosen on the basis of the previous calculation by Chaudhuri et al.\textsuperscript{38} We found that when an additional O adatom (O3, red) is adsorbed far from the ketone group, the wear mechanism and the associated energy behave similar to what observed when only two O atoms are adsorbed (Figure 3), i.e., a carbon chain is formed upon the displacement of the O atom forming a double bond with the first detached carbon. On the contrary, when an additional O adatom (O4, blue) forms a bond with a first neighboring site of the ketone group, the detachment of the carbon chain is inhibited and only the C atom belonging to the ketone group is detached through the formation of a CO molecule (Figure 4), as observed in the case of the C(001) surface. This could be because the increase in the density of adsorbed oxygen follows by the reduction of bond strength of C−C bonds involving the O adsorption. The C−C(−O) bond is broken to create a CO molecule rather than a carbon chain. As a result, the energy shows a different trend with respect to that shown in Figure 3. The result suggested that the wear mechanism as well as the length of the C chain is affected by the adsorption of other species around the ketone group.

We further extend the study of the wear mechanism on the C(110) surface by considering the presence of different chemical species and dissociated molecules including OH, OOH, H\textsubscript{2}, and H\textsubscript{2}O, as shown in Figure 5a. These species are highly present in the working environment as products of the dissociative adsorption of atmospheric gases or precursor molecules used for diamond growth.\textsuperscript{39−41} Starting with the most vulnerable configuration to wear (ketone configuration in Figure 1), we straightforwardly substitute the oxygen atom with different species and atoms. We found that when the \pdfpage{5}oxygen atom in the ketone group is replaced by a H atom or OH group, wear cannot be initiated by the out-of-plane displacement of these groups: the displaced atoms are detached from the surface without causing any deformation for the C−C bonds of the substrate. The energy (Figure 5b) and force (Figure S3) trends are similar for H/OH. This could be because these groups form single bonds with the substrate. These single bonds are weaker than the O═C double bond presented in the ketone configuration, so the C atom dragged by the pulled atom cannot withstand the restoring force from the substrate. Owing to the analogous zigzag arrangements of carbon atoms atop both the C(110) and C(111) surfaces, a comparable wear mechanism is anticipated to emerge for these surfaces upon interactions with different adsorption species. On the other hand, we believe that the wear mechanism on the C(001) surface, characterized by the formation of CO molecules, is unlikely to undergo significant alteration upon the adsorption of alternative species. This is due to the fact that (1) the carbon atoms on top of C(001) surface do not adopt a connected zigzag chain configuration. Consequently, they are prone to dissociation upon application of an upward force. (2) The adsorption of other species will weaken the C═C double bonds. This outcome results from the necessity for carbon atoms to share electrons to form a covalent bond with other species (H or OH). As a result, the bond strength is weakened and will facilitate more the formation of CO molecule. Second, if the O forms a double bond with the C atom, it could reduce the strength of other C−C bonds as more electrons are shared to form the O═C double bond. As shown in Figure 5b (black line), the energy cost necessary to pull out of the surface the carbon chain terminating with an O atom (resulting, e.g., from the dissociation of a OH group) is much lower than the energy cost to detach the OH or H groups from the surface. Thus, the pulling of H or OH groups is not possible to promote the detachment of any C atom from the surface. Therefore, we conclude that the dragging specie to initiate wear should form a double bond with the displaced carbon atom, as stronger than the restoring C−C bonds.

%FIG6
\begin{center}
\includegraphics[width=0.95\linewidth,height=0.55\textheight,keepaspectratio]{figures/fig6.png}
\end{center}
\noindent\textbf{Figure 6.} Different initial adsorption configurations of a silicate cluster on the C(110) surface (a). Wear mechanism caused by the detachment of the silicate cluster on the C(110) surface (b). C (gray), O (red), H (white), and Si (yellow).
\decompnote{Figure 6 is printed across the top of p.\,5, above the first text line of that page; it is placed after the paragraph that runs from p.\,4 onto p.\,5.}
\alttext{(a) Six ball-and-stick panels in light-blue frames, columns ``Case \#1'', ``Case \#2'', ``Case \#3'': top views (upper row) and side views (lower row) of a silicate cluster (Si yellow, O red, H white) on C(110) (C black/grey). Case \#1: the Si is bonded through one O to a C atom that stands out of the surface; Case \#2: two O atoms of the cluster are bonded to the same out-of-plane C atom; Case \#3: the Si sits close to the surface and is drawn with two bonds down to it. Each side view also shows an O atom on the surface to the right.
(b) Line plot, x ``Displacement (Å)'' 0 to 2.5, y ``Energy (eV)'' −1 to 7; legend Case \#1 blue squares, Case \#2 red circles (semi-transparent), Case \#3 grey triangles on a black line. Case \#1: 0, 0.10, 0.34, 1.58, 2.55 at 1.0 Å, then about 2.5–2.6 eV to 2.0 Å, where it ends. Case \#2: 0, 0.27, 0.95, 1.84, 2.21, 2.77, 3.47 at 1.5 Å, 3.80, 3.98, 4.06 and 4.06 eV at 2.5 Å. Case \#3: 0, 0.26, 0.99, 2.29 at 0.75, 2.45, 3.19 at 1.25 Å (local maximum), 2.78, 2.87, 3.08, 3.53 and 4.26 eV at 2.5 Å. Dotted vertical lines with call-outs 1 (0 Å), 2 (0.75 Å), 3 (1.25 Å), 4 (1.5 Å) and 5 (2.5 Å) connect to colour-framed side-view insets: initial structures at the left; at 0.75 and 1.5 Å blue, red and black frames with the cluster lifted (in some the C atom it holds lifted too); at 1.25 Å (bottom, red frame) the cluster on top of a short vertical C chain; at 2.5 Å a red frame with the cluster at the top of a longer C chain and a black frame with the cluster detached above the surface.}
\figuredata{b}{figures/fig6b.csv}{figures/fig6b.py}{tier 3, digitized from the embedded 300 ppi raster image (no vector data in the PDF); marker centres located by colour masks and axis fits on the major ticks; covered markers read from their visible edges or from the curve on both sides (flagged in the CSV); see figures/fig6b\_extract.py}

In summary, the formation of wear and its mechanisms of diamond depend greatly on surface orientations and adsorbed species. Among the three surfaces studied, the C(110) surface was identified as the most susceptible surface to wear at both the clean and the O-adsorbed structures. Under oxygen adsorption, only the ketone configuration can initiate wear on diamond surfaces through CO desorption or carbon chain. In addition, the contribution of adsorbed species such as H/OH/OOH, and oxygen coadsorption on mechanism of wear was also investigated. Although H/OH/OOH adsorption shows a minor effect on diamond wear, the coadsorption of oxygen atoms can alter the wear mechanism induced by ketone configuration from carbon chain to CO desorption. This is the synergistic effect from multiple C−O bond formation, resulting in a reduction of C−C bond strength and facilitating C−C bond rupture. The calculation results highlight the intricate interplay between surface properties and adsorbed species on the wear mechanisms of diamond.

\noindent\textbf{3.2. Wear Due to the Interaction with a Silicate.} Experimental studies have indicated that diamond can be worn when the sliding against silicate.\textsuperscript{5,25,42} The initiation of wear is associated with the formation of chemical bonds across the interface.\textsuperscript{9,36} In such a case, the local environment where Si−O−C bond formation occurs becomes extremely useful to examine atomistic mechanisms of wear. During sliding, the relative movement of silicate against diamond leads to the stretching of the Si−O−C bonds to promote wear.\textsuperscript{24} The impact of the counter surface sliding thus can be directed to the movement of Si atom lifting up the oxygen atom through the Si−O−C connection. To investigate this situation, we adopt a simplified model containing a silicate cluster interacting with the C(110) surface, which was found to be the most easily wearable. We consider three different adsorption configurations of the silicate cluster shown in Figure 6. In all of the initial configurations, the silicate cluster is positioned on top of an out-of-plane carbon atom. As shown in Figure 6b, when the silicate cluster bonds to the surface through single Si−O−C bonds (one in case \#1 in blue and two in case \#3 in black), the Si−O bonds broken after the Si displacement, leading to a sudden drop of force (Figure S4).

On the other hand, when the silicate cluster is bonded with the surface through two Si−O−C single bonds involving the same C atom, known as bidentate structure, the pulling up of the silicate fragment can lead to the formation of a C chain as shown in Figure 6b (case \#2 in red). This mechanism is similar to that observed when the O atom of the ketone O═C group \pdfpage{6}is pulled, as described in Section 3.1, and we observe a similar trend in the energy (Figure 6b) and force (Figure S4) curves. Particularly, the force on the Si atom peaks each time the carbon chain is stressed and drops suddenly when the C−C underlying bond is broken, and the longer the carbon chain, the longer the displacement that it takes for the force to rise. The result suggested that the two single bonds involving the same C atom can pull strongly enough to cause its detachment from the surface, as observed in the C═O double bond.

%FIG7
\begin{center}
\includegraphics[width=0.95\linewidth,height=0.55\textheight,keepaspectratio]{figures/fig7.png}
\end{center}
\noindent\textbf{Figure 7.} Wear mechanism caused by the adsorption of oxygen on the C(110) surface under different stress levels.
\decompnote{Figure 7 is printed across the top of p.\,6, above the first text line of that page; it is placed after the paragraph that runs from p.\,5 onto p.\,6.}
\alttext{Two line plots, each x ``Displacement (Å)'' 0 to 2.5 and y ``Energy (eV)''; legend 0\% grey circles (black line), 3\% red squares, 7\% blue triangles, 10\% green downward triangles; across the top of each a row of four side-view insets framed black, red, blue and green and labelled 0\%, 3\%, 7\%, 10\% (C black/grey, O red), and a blue-framed inset labelled ``initial'' (the ketone-configured surface with an O bridging a bond).
(a) ``(a) Compression'', y −6 to 7. 0\%: 0, 0.33, 1.48, 2.0, 2.63, about 2.45 from 1.25 to 1.75 Å, 2.70, 3.28, 3.05 eV. 3\%: 0.62, 1.62, 2.11, 2.82 at 1.0 Å, a drop to 0.95 at 1.25, 1.16, 1.92, 2.98, 3.31 and 2.61 eV. 7\%: 1.18 at 0.25, −0.18 at 0.5, 0.92, 0.45, 0.57, 0.78, 1.57, 2.83 at 2.0 Å, 1.78 and 0.34 eV. 10\%: entirely negative after 0: −1.96, −1.87, −1.52, a minimum of −4.96 eV at 1.0 Å, −4.86, −4.54, −3.66, −2.32, −2.45, −2.44 eV. The four end-state insets all show a chain of C atoms with O at the top rising from the surface (inclined for 3\%).
(b) ``(b) Expansion'', y −1 to 7. 0\%: 0, about 0.35, 1.47, 1.99, 2.63, about 2.45 from 1.25 to 1.75 Å, 2.69, 3.24, 3.06 eV. 3\%: 0.64, 1.63, 2.12, 2.82, 0.95 at 1.25 Å, 1.14, 1.87, 3.00, 3.32, 2.60 eV. 7\%: 0.13, 0.79, 1.94, 3.29, then a plateau of about 3.9–3.97 eV from 1.25 Å. 10\%: 0.23, 0.67, 1.59, 3.03, a maximum of 4.06 eV at 1.25 Å, then about 3.76–3.84 eV. Insets: 0\% and 3\% show a vertical chain of C atoms with O on top; 7\% and 10\% show a detached C−O pair above the surface.}
\figuredata{a}{figures/fig7a.csv}{figures/fig7a.py}{tier 3, digitized from the embedded 300 ppi raster image (no vector data in the PDF); marker centres located by colour masks and axis fits on the major ticks; covered markers read from their visible edges or from the curve on both sides (flagged in the CSV); see figures/fig7a\_extract.py}
\figuredata{b}{figures/fig7b.csv}{figures/fig7b.py}{tier 3, digitized from the embedded 300 ppi raster image (no vector data in the PDF); marker centres located by colour masks and axis fits on the major ticks; covered markers read from their visible edges or from the curve on both sides (flagged in the CSV); see figures/fig7b\_extract.py}

In comparison with the O-adsorbed C(110) surface, the single Si−O bond plays a role similar to that of the ether configuration while the bidentate structure can be a resemblance of the ketone configuration. In this context, the wear mechanisms are deeply governed by the bond strength of the bond that the pulled atoms form with carbon. Thus, the interactions of single and double Si−O bonds with the carbon atoms belonging to the other surfaces of diamond C(001) and C(111) can result in similar effects to those described for the ether and ketone adsorption configurations of oxygen on these two surfaces.

\noindent\textbf{3.3. Effect of Stress on the Wear Mechanism.} In tribological conditions, the surface can be subjected to stress as a result of sliding, loads, and preparation of the material,\textsuperscript{23,37,43,44} which can lead to the peeling-off of carbon atoms from the diamond surface.\textsuperscript{45} Consequently, it is crucial to investigate the impact of stress on the wear mechanism of diamond surfaces. In the present study, stress is introduced by varying the lateral size of the supercell with four stress levels considered, including 0\%, 3\%, 7\%, and 10\%. Such high values of stress can be experienced at sharp asperity−asperity contacts where the local load and temperature are extremely high, and deformation takes place.\textsuperscript{46} Furthermore, the application of extreme stresses, especially 7\% and 10\%, aims to accelerate significant changes within the system. The objective is to understand the influence of stress on the atomistic wear mechanisms of diamond. The evolution of energy and diamond structures subjected to stress are presented in Figure 7.

In the case of compression (as depicted in Figure 7a and the forces in Figure S5), our findings indicate that the wear mechanism remains unchanged with the detachment of oxygen atoms leading to the formation of long carbon chains. However, the stress can significantly alter the energy cost of wear. Specifically, the energy required to detach the oxygen atom decreases as the surface is compressed, particularly under high-stress conditions (7\% and 10\%). Interestingly, when the surface is compressed at 10\%, it becomes energetically favorable for carbon atoms to detach from the surface, as shown in Figure 7a. Therefore, our result suggests that under extremely high load, the wear on the diamond surface can become severe.

On the other hand, the wear mechanism can undergo significant alteration when the surface is subjected to expansion. As depicted in Figure 7b, the formation of long C chains remains the predominant wear mechanism at lower stress levels (0\% and 3\%). However, when the stress level is increased, it fosters the formation of CO molecules rather than long C chains. Our findings reveal that both energy and force are augmented as the wear mechanism involves the formation of CO molecules. Moreover, there is a negligible difference observed between stress levels of 7\% and 10\%. Compared with the compression, the expansion promotes the formation of CO molecules rather than carbon chains. This is because the C−C bond distance is shortened under compression and elongated under expansion. The bond elongation makes the C−C bond weaker and thus unable to withstand the pulling forces as in the compressed case. Thus, the C−C bond at the surface is detached, and a CO molecule is produced. Furthermore, the expansion leads to the reduction of the diamond density, making it unlikely to release its atoms as a chain, as found for the compression.

\section*{4. CONCLUSIONS}

In conclusion, the atomistic mechanisms of diamond wear are investigated by first-principles calculations considering different surface orientations, adsorbates, and stress levels in the carbon film. Our findings can be summarized as follows.
\begin{itemize}
\item The primary mechanism of wear involves the detachment of carbon chains from the surface. The C(110) surface is the most easily worn. Meanwhile the Pandey-reconstructed C(111) surface shows the best resistance to wear.
\item The restoring force on the displaced C atom gradually increases until a C−C bond breaks on the diamond \pdfpage{7}surface. A carbon chain is then formed by repeating this mechanism. Interestingly, the elastic properties of the formed chain resemble those of springs added in series: the longer the carbon chain, the lower the restoring force and energy increase.
\item In the presence of O, H, and OH adsorbates, the formation of a chain pulled by a ketone, C═O, requires lower energy than on the clean diamond surfaces. On the contrary, the single adsorbate−C bonds turn out to be less strong than the surface C−C bonds; thus, only the pulled adsorbate is detached, without causing any distortion of the underlying carbon surface.
\item When the adsorbate is a silicate, our results indicate that single Si−O−C or Si−O bonds are not enough to break C−C bonds on the diamond surface. Two Si−O−C bonds involving the same C atom in a bidentate structure are instead able to promote the C detachment from the surface and initiate the wear process through chain formation.
\end{itemize}

The stress level present on the diamond surface can significantly impact its wear resistance. A compressive stress decreases the energy cost for the detachment of carbon atoms, making wear an energetically favorable process under highly stressed conditions. Conversly, when the surface is subjected to expansion, the wear mechanism of carbon chain formation is inhibited, and removal of C atoms occurs through the formation of CO molecules at a high energy cost.
\decompnote{``Conversly'' is the paper's spelling.}

\section*{■ ASSOCIATED CONTENT}

\subsection*{Supporting Information}

The Supporting Information is available free of charge at https://pubs.acs.org/doi/10.1021/acs.langmuir.3c01800.

\begin{quote}
Evolution of forces during the pulling processes of clean and O-adsorbed diamond (110), (001), and (111) surfaces, force changes during the detachment of O/H/OH, and a silicate cluster on C(110) surface, force changes by the adsorption of oxygen on C(110) surface under different stress levels (PDF)
\end{quote}
\decompnote{The Supporting Information PDF (Figures S1–S5, the force curves) is not in the local collection; see README.}

\section*{■ AUTHOR INFORMATION}

\subsection*{Corresponding Authors}

\noindent\textbf{M. C. Righi} − \textit{Department of Physics and Astronomy, University of Bologna, 40127 Bologna, Italy;} orcid.org/0000-0001-5115-5801; Email: clelia.righi@unibo.it

\noindent\textbf{Nam V. Tran} − \textit{School of Material Science and Engineering, Nanyang Technological University, 639798, Singapore;} orcid.org/0000-0002-1756-1475; Email: vannam.tran@ntu.edu.sg

\subsection*{Author}

\noindent\textbf{Huong T. T. Ta} − \textit{Department of Physics and Astronomy, University of Bologna, 40127 Bologna, Italy;} orcid.org/0000-0002-2675-0681

\noindent Complete contact information is available at: https://pubs.acs.org/10.1021/acs.langmuir.3c01800

\subsection*{Notes}

The authors declare no competing financial interest.

\section*{■ ACKNOWLEDGMENTS}

These results are part of the “Advancing Solid Interface and Lubricants by First-Principles Material Design (SLIDE)” project that has received funding from the European Research Council (ERC) under the European Union’s Horizon 2020 research and innovation program (Grant Agreement No. 865633). PRACE is also acknowledged for computational resources.

\section*{■ REFERENCES}

\begin{hangpars}\small
(1) Hainsworth, S. V.; Uhure, N. J. Diamond like Carbon Coatings for Tribology: Production Techniques, Characterisation Methods and Applications. \textit{Inter. Mater. Rev.} \textbf{2007}, \textit{52}, 153−174.

(2) Wheeler, D. W. Applications of Diamond to Improve Tribological Performance in the Oil and Gas Industry. \textit{Lubricants} \textbf{2018}, \textit{6}, 84.

(3) Erdemir, A.; Donnet, C. Tribology of Diamond-like Carbon Films: Recent Progress and Future Prospects. \textit{J. Phys. D: Appl. Phys.} \textbf{2006}, \textit{39}, R311−R327.

(4) Liu, Y.; Erdemir, A.; Meletis, E. I. A Study of the Wear Mechanism of Diamond-like Carbon Films. \textit{Surf. Coat. Technol.} \textbf{1996}, \textit{82}, 48−56.

(5) Khurshudov, A. G.; Kato, K.; Koide, H. Wear of the AFM Diamond Tip Sliding against Silicon. \textit{Wear} \textbf{1997}, \textit{203}−\textit{204}, 22−27.

(6) Grillo, S. E.; Field, J. E.; Van Bouwelen, F. M. Diamond Polishing: The Dependency of Friction and Wear on Load and Crystal Orientation. \textit{J. Phys. D: Appl. Phys.} \textbf{2000}, \textit{33}, 985−990.

(7) Yan, J.; Syoji, K.; Tamaki, J. Some Observations on the Wear of Diamond Tools in Ultra-Precision Cutting of Single-Crystal Silicon. \textit{Wear} \textbf{2003}, \textit{255} (7−12), 1380−1387.

(8) Wang, L.; Cui, L.; Lu, Z.; Zhou, H. Understanding the Unusual Friction Behavior of Hydrogen-Free Diamond-like Carbon Films in Oxygen Atmosphere by First-Principles Calculations. \textit{Carbon} \textbf{2016}, \textit{100}, 556−563.

(9) Cutini, M.; Forghieri, G.; Ferrario, M.; Righi, M. C. Adhesion, Friction and Tribochemical Reactions at the Diamond-Silica Interface. \textit{Carbon} \textbf{2023}, \textit{203}, 601−610.

(10) Zilibotti, G.; Righi, M. C.; Ferrario, M. Ab Initio Study on the Surface Chemistry and Nanotribological Properties of Passivated Diamond Surfaces. \textit{Phys. Rev. B} \textbf{2009}, \textit{79} (7), 075420.

(11) Erdemir, A. The Role of Hydrogen in Tribological Properties of Diamond-like Carbon Films. \textit{Surf. Coat. Technol.} \textbf{2001}, \textit{146}−\textit{147}, 292−297.

(12) Tran, N. V.; Righi, M. C. Ab Initio Insights into the Interaction Mechanisms between H2, H\textsubscript{2}O, and O2Molecules with Diamond Surfaces. \textit{Carbon} \textbf{2022}, \textit{199}, 497−507.

(13) Enriquez, J. I.; Muttaqien, F.; Michiuchi, M.; Inagaki, K.; Geshi, M.; Hamada, I.; Morikawa, Y. Oxidative Etching Mechanism of the Diamond (100) Surface. \textit{Carbon} \textbf{2021}, \textit{174}, 36−51.

(14) John, P.; Polwart, N.; Troupe, C. E.; Wilson, J. I. B. The Oxidation of Diamond: The Geometry and Stretching Frequency of Carbonyl on the (100) Surface. \textit{J. Am. Chem. Soc.} \textbf{2003}, \textit{125} (22), 6600−6601.

(15) Shirafuji, J.; Sugino, T. Electrical Properties of Diamond Surfaces. \textit{Diam. Related Mater.} \textbf{1996}, \textit{5}, 706−713.

(16) Zheng, J.; Xie, X.; Wee, A.; Ping Loh, K. Oxygen-Induced Surface State on Diamond 100. \textit{Diam. Related Mater.} \textbf{2001}, \textit{10}, 500−505.

(17) Loh, K. P.; Xie, X. N.; Yang, S. W.; Zheng, J. C. Oxygen Adsorption on (111)-Oriented Diamond: A Study with Ultraviolet Photoelectron Spectroscopy, Temperature-Programmed Desorption, and Periodic Density Functional Theory. \textit{J. Phys. Chem. B} \textbf{2002}, \textit{106} (20), 5230−5240.

(18) Guo, X.; Yuan, S.; Wang, X.; Jin, Z.; Kang, R. Atomistic Mechanisms of Chemical Mechanical Polishing of Diamond (100) in Aqueous H\textsubscript{2}O\textsubscript{2}/Pure H\textsubscript{2}O: Molecular Dynamics Simulations Using Reactive Force Field (ReaxFF). \textit{Comput. Mater. Sci.} \textbf{2019}, \textit{157}, 99−106.

(19) Rani, R.; Panda, K.; Kumar, N.; Kozakov, A. T.; Kolesnikov, V. I.; Sidashov, A. V.; Lin, I. N. Tribological Properties of \pdfpage{8}Ultrananocrystalline Diamond Films: Mechanochemical Transformation of Sliding Interfaces. \textit{Sci. Rep.} \textbf{2018}, \textit{8}, 283.

(20) Pastewka, L.; Moser, S.; Gumbsch, P.; Moseler, M. Anisotropic Mechanical Amorphization Drives Wear in Diamond. \textit{Nat. Mater.} \textbf{2011}, \textit{10} (1), 34−38.

(21) Kawaguchi, K.; Wang, Y.; Xu, J.; Ootani, Y.; Higuchi, Y.; Ozawa, N.; Kubo, M. Atom-by-Atom and Sheet-by-Sheet Chemical Mechanical Polishing of Diamond Assisted by OH Radicals: A Tight-Binding Quantum Chemical Molecular Dynamics Simulation Study. \textit{ACS Appl. Mater. Interfaces} \textbf{2021}, \textit{13} (34), 41231−41237.

(22) Huisman, W. J.; Peters, J. F.; De Vries, S. A.; Vlieg, E.; Yang, W.-S.; Derry, T. E.; Van Der Veen, J. F. Structure and Morphology of the As-Polished Diamond (111)-1 × 1 Surface. \textit{Surf. Sci.} \textbf{1997}, \textit{387}, 342−353.

(23) Liu, J.; Jiang, Y.; Grierson, D. S.; Sridharan, K.; Shao, Y.; Jacobs, T. D. B.; Falk, M. L.; Carpick, R. W.; Turner, K. T. Tribochemical Wear of Diamond-Like Carbon-Coated Atomic Force Microscope Tips. \textit{ACS Appl. Mater. Interfaces} \textbf{2017}, \textit{9} (40), 35341−35348.

(24) Peguiron, A.; Moras, G.; Walter, M.; Uetsuka, H.; Pastewka, L.; Moseler, M. Activation and Mechanochemical Breaking of C-C Bonds Initiate Wear of Diamond (110) Surfaces in Contact with Silica. \textit{Carbon} \textbf{2016}, \textit{98}, 474−483.

(25) Thomas, E. L. H.; Mandal, S.; Brousseau, E. B.; Williams, O. A. Silica Based Polishing of \{100\} and \{111\} Single Crystal Diamond. \textit{Sci. Technol. Adv. Mater.} \textbf{2014}, \textit{15}, 035013.

(26) Giannozzi, P.; et al. QUANTUM ESPRESSO: A Modular and Open-Source Software Project for Quantum Simulations of Materials. \textit{J. Phys. Condens. Mater.} \textbf{2009}, \textit{21} (39), 395502.

(27) Perdew, J. P.; Burke, K.; Ernzerhof, M. Generalized Gradient Approximation Made Simple. \textit{Phys. Rev. Lett.} \textbf{1996}, \textit{77}, 3865.

(28) Bučko, T.; Hafner, J.; Lebègue, S.; Ángyán, J. G. Improved Description of the Structure of Molecular and Layered Crystals: Ab Initio DFT Calculations with van Der Waals Corrections. \textit{J. Phys. Chem. A} \textbf{2010}, \textit{114} (43), 11814−11824.

(29) Grimme, S. Semiempirical GGA-Type Density Functional Constructed with a Long-Range Dispersion Correction. \textit{J. Comput. Chem.} \textbf{2006}, \textit{27} (15), 1787−1799.

(30) Lebedeva, I. V.; Lebedev, A. V.; Popov, A. M.; Knizhnik, A. A. Comparison of Performance of van Der Waals-Corrected Exchange-Correlation Functionals for Interlayer Interaction in Graphene and Hexagonal Boron Nitride. \textit{Comput. Mater. Sci.} \textbf{2017}, \textit{128}, 45−58.

(31) Xiang, Y.; Jiang, M.; Xiao, H.; Xing, K.; Peng, X.; Zhang, S.; Qi, D. C. A DFT Study of the Surface Charge Transfer Doping of Diamond by Chromium Trioxide. \textit{Appl. Surf. Sci.} \textbf{2019}, \textit{496}, 143604.

(32) Ayestarán Latorre, C.; Ewen, J. P.; Dini, D.; Righi, M. C. Ab Initio Insights into the Interaction Mechanisms between Boron, Nitrogen and Oxygen Doped Diamond Surfaces and Water Molecules. \textit{Carbon} \textbf{2021}, \textit{171}, 575−584.

(33) Rivero, P.; Shelton, W.; Meunier, V. Surface Properties of Hydrogenated Diamond in the Presence of Adsorbates: A Hybrid Functional DFT Study. \textit{Carbon} \textbf{2016}, \textit{110}, 469−479.

(34) De La Pierre, M.; Bruno, M.; Manfredotti, C.; Nestola, F.; Prencipe, M.; Manfredotti, C. The (100), (111) and (110) Surfaces of Diamond: An Ab Initio B3LYP Study. \textit{Mol. Phys.} \textbf{2014}, \textit{112} (7), 1030−1039.

(35) Saito, M.; Miyamoto, Y.; Oshiyama, A. Stability and Lattice Vibrations of the Clean Diamond (111) Surface. \textit{Surf. Sci.} \textbf{1999}, \textit{427}−\textit{428}, 53−57.

(36) Milne, Z. B.; Schall, J. D.; Jacobs, T. D. B.; Harrison, J. A.; Carpick, R. W. Covalent Bonding and Atomic-Level Plasticity Increase Adhesion in Silicon-Diamond Nanocontacts. \textit{ACS Appl. Mater. Interfaces} \textbf{2019}, \textit{11} (43), 40734−40748.

(37) Gotsmann, B.; Lantz, M. A. Atomistic Wear in a Single Asperity Sliding Contact. \textit{Phys. Rev. Lett.} \textbf{2008}, \textit{101} (12), 125501.

(38) Chaudhuri, S.; Hall, S. J.; Klein, B. P.; Walker, M.; Logsdail, A. J.; Macpherson, J. V.; Maurer, R. J. Coexistence of Carbonyl and Ether Groups on Oxygen-Terminated (110)-Oriented Diamond Surfaces. \textit{Commun. Mater.} \textbf{2022}, \textit{3} (1), 6.

(39) Konicek, A. R.; Grierson, D. S.; Sumant, A. V.; Friedmann, T. A.; Sullivan, J. P.; Gilbert, P. U. P. A.; Sawyer, W. G.; Carpick, R. W. Influence of Surface Passivation on the Friction and Wear Behavior of Ultrananocrystalline Diamond and Tetrahedral Amorphous Carbon Thin Films. \textit{Phys. Rev. B} \textbf{2012}, \textit{85} (15), 155448.

(40) Park, S. J.; Lee, K. R.; Ko, D. H. Tribochemical Reaction of Hydrogenated Diamond-like Carbon Films: A Clue to Understand the Environmental Dependence. \textit{Tribol. Int.} \textbf{2004}, \textit{37} (11−12), 913−921.

(41) Uchidate, M.; Liu, H.; Iwabuchi, A.; Yamamoto, K. Effects of Water Environment on Tribological Properties of DLC Rubbed against Stainless Steel. \textit{Wear} \textbf{2007}, \textit{263} (7−12), 1335−1340.

(42) Kumar, A.; Melkote, S. N. Wear of Diamond in Scribing of Multi-Crystalline Silicon. \textit{J. Appl. Phys.} \textbf{2018}, \textit{124} (6), 5037106.

(43) Shu, K.; Zhang, C.; Sun, D.; Wu, B.; Li, Z.; Wu, J.; Zheng, D.; Gu, L. Contact Cracking and Delamination of Thin Solid Lubricating Films on Deformable Substrates. \textit{Tribol. Int.} \textbf{2021}, \textit{163}, 107181.

(44) Almeida, F. A.; Carrapichano, J. M.; Fernandes, A. J. S.; Sacramento, J.; Silva, R. F.; Oliveira, F. J. Nanocrystalline CVD Diamond Coatings for Drilling of WC-Co Parts. \textit{Int. J. Refract. Metals Hard Mater.} \textbf{2011}, \textit{29} (5), 618−622.

(45) Mckenzie, D. R.; Muller, D.; Pailthorpe, B. A. Compressive-Stress-Induced Formation of Thin-Film Tetrahedral Amorphous Carbon. \textit{Phys. Rev. Lett.} \textbf{1991}, \textit{67}, 773.

(46) Bay, N.; Wanheim, T. Real Area of Contact and Friction Stress at High Pressure Sliding Contact. \textit{Wear} \textbf{1976}, \textit{38}, 201−209.
\end{hangpars}
\decompnote{Ref.\,19: the word is split at a line end on p.\,7 as ``Ultra-/nanocrystalline''; it is written without the hyphen, as ``Ultrananocrystalline'' is printed in ref.\,39. Refs.\,21, 30 and 45 also break at a hyphen (``Tight-/Binding'', ``Exchange-/Correlation'', ``Compressive-/Stress''); those hyphens are kept, as the hyphenated compound reads as part of the title.}

\end{document}
